\documentclass[margin=10pt]{standalone}
\usepackage[utf8]{inputenc}
\usepackage{pgfplots}
\usepackage{siunitx} % Per formattare le unità di misura

\pgfplotsset{compat=1.18}

\begin{document}
	
	\begin{tikzpicture}
		\begin{axis}[
			width=12cm,
			height=8cm,
			grid=major,
			xlabel={Tempo (\unit{\second})},
			ylabel={Ampiezza (\unit{\volt})},
			title={\textbf{Particelle alpha 02}},
			% Trasformiamo i valori di tempo in microsecondi per leggibilità
			% scaled x ticks=false,
			% x tick label style={/pgf/number format/fixed},
			legend pos=north east,
			no markers, % Rimuove i pallini, lascia solo la linea
			]
			
			\addplot [
			blue, 
			thick,
			% Poiché il file ha 10.000 punti, leggiamo un punto ogni 10
			% per non saturare la memoria di LaTeX (rimuovi se vuoi tutti i punti)
			each nth point={50}, 
			filter discard warning=false
			] table [
			x=TIME, 
			y=CH1, 
			col sep=comma
			] {tek0004CH1.csv};
			
			\legend{Segnale CH1}
			
		\end{axis}
	\end{tikzpicture}
	
\end{document}